\section{实验结果与应用}

\subsection{语言建模性能}

\subsubsection{LM1B 数据集}

在 LM1B 数据集（上下文长度 128）上，MDLM 显著优于所有先前的离散扩散模型基线（包括最接近的基线 SEDD），并\textbf{接近了 AR 模型的困惑度}。

\begin{knowledgebox}{为什么困惑度重要}
困惑度（Perplexity）是语言建模中最核心的评估指标。经验表明，如果一个模型在训练/验证集上取得了好的困惑度分数，它在下游任务（如编程、推理）上的表现也会更好。因此，困惑度是衡量语言模型质量的可靠代理指标。
\end{knowledgebox}

当上下文长度增加到 1024 时，MDLM 仍然优于 SEDD，但与 AR 模型之间存在一定差距。不过，考虑到早期的扩散语言模型与 AR 基线差距巨大，这一结果已经非常令人鼓舞。

\subsubsection{零样本迁移能力}

将模型在 OpenWebText 上训练后，在 PTB、arXiv、PubMed 等 7 个未见过的数据集上评估零样本似然：
\begin{itemize}
    \item MDLM 在所有数据集上一致优于 SEDD
    \item 在 7 个数据集中的 3 个上，MDLM 的表现\textbf{优于 AR 基线}
\end{itemize}

\subsection{大规模验证：LaDa（8B MDLM）}

LaDa（Large Language Diffusion Models）是 MDLM 在 80 亿参数规模上的验证。它本质上就是一个 8B 的 MDLM 模型。

\begin{importantbox}{MDLM 在大规模上的表现}
\begin{itemize}
    \item \textbf{多任务语言理解}（MMLU）：MDLM 在整个训练过程中\textbf{一致超过 AR 基线}
    \item \textbf{推理任务}：MDLM 的性能扩展曲线呈\textbf{指数增长}，而 AR 基线仅为线性增长
    \item \textbf{常识问答和编程}：MDLM 与 AR 基线\textbf{持平}或\textbf{略低}
\end{itemize}
尽管在困惑度上仍有差距，但在实际下游任务上，MDLM 在多个方面超越了 AR 模型。
\end{importantbox}

\subsection{超越语言：离散数据建模}

\begin{knowledgebox}{左右偏置与数据结构}
自然语言具有固有的从左到右的偏置——这是人类说话和写作的方式。因此 AR 模型在语言任务上有天然优势。但在蛋白质生成、药物发现、分子生成等领域，数据\textbf{不存在左右偏置}，这正是 MDLM 大放异彩的场景。
\end{knowledgebox}

在分子生成任务中（GenMol 论文），基于 MDLM 的模型展示了一个显著优势：通过调节采样温度，可以在分子\textbf{质量}和\textbf{多样性}之间获得一条 Pareto 前沿。相比之下，AR 基线只能给出一个固定的质量-多样性点，无法进行这种灵活的权衡。

\subsection{可控生成}

MDLM 天然适合可控生成。原因在于：

\begin{enumerate}
    \item 在去噪过程中，模型会填充所有 mask 位置，产生对最终样本的\textbf{完整预测}
    \item 由此可以获得关于最终样本的高层信息
    \item 利用这些信息，可以在采样过程中引导生成具有特定属性的样本
\end{enumerate}

\begin{warningbox}{AR 模型难以进行可控采样}
AR 模型在可控生成方面存在众所周知的困难。因为在逐词生成的过程中，模型无法预见未来的全局结构，因此很难在生成过程中施加全局约束。MDLM 的双向注意力机制和全局去噪预测使其在可控生成上具有本质优势。
\end{warningbox}

\subsection{本章小结}

实验结果表明：（1）MDLM 在小规模上接近 AR 性能，在大规模下游任务（特别是推理）上甚至超过 AR；（2）在无左右偏置的离散数据（如分子生成）上，MDLM 表现尤为突出；（3）MDLM 的采样灵活性使其适合可控生成等应用场景。

\newpage
